% !PRINTED 94
\weberpage{94}
Es hat dann, wie man aus
\[
\frac{\omega'_2}{\omega'_1} = \frac{c\omega_1+\partial\omega_2}{a\omega_1+b\omega_2}
\]
durch Trennung des reellen vom imaginären Teil erkennt, der
imaginäre Teil von $\omega_2:\omega_1$ dasselbe Vorzeichen, wie der von
$\omega'_2:\omega'_1$ (das positive).

Setzen wir
\weq{7}{\begin{gathered}
a_1 = \partial a'_1 - b a'_2 \\
a_2 = -c a'_1 + a a'_2
\end{gathered}}
\weq{8}{\begin{gathered}
b_1 = \partial b'_1 - b b'_2 - m'b\partial \\
b_2 = -cb'_1 + ab'_2 - m'ac,
\end{gathered}}
so schließt man aus (3), daß die Funktion $T(u)$ nicht nur den
Bedingungen (1), sondern auch den aus (1) durch Vertauschung
von $\omega'_1$, $\omega'_2$, $a'_1$, $a'_2$, $b'_1$, $b'_2$ mit $\omega_1$, $\omega_2$, $a_1$, $a_2$, $b_1$, $b_2$ hervorgehenden
Gleichungen, d.\,h. den Gleichungen (1), § 17, genügt. Sie ist
also gleichzeitig eine $T$-Funktion der Perioden $\omega'_1$, $\omega'_2$ und der
Perioden $\omega_1$, $\omega_2$, was wir durch folgende Gleichung andeuten:
\weq{9}{T'(u,\,\omega'_1\,\omega'_2) = T(u,\,\omega_1,\,\omega_2).}
% [sic] Printed without a comma between the two primed periods.
Es ist aber nach (5) und (7)
\[
a_2\omega_1 - a_1\omega_2 = (a'_2\omega'_1 - a'_1\omega'_2)(a\partial - bc),
\]
also, wenn $m$ die Ordnung von $T$ ist,
\weq{10}{m = m'n.}
Nach (8) ist die Charakteristik $(g_1, g_2)$ von $T$, wenn $(g'_1, g'_2)$
die von $T'$ ist (§ 17),
\weq{11}{(g_1, g_2) = (\partial g'_1 - b g'_2 - m'b\partial,\ -c g'_1 + a g'_2 - m'ac).}
\gesperrt{Unter der Transformation der $T$-Funktionen versteht
man die Darstellung der Funktionen $T'$ mit den Perioden
$\omega'_1$, $\omega'_2$ durch $T$-Funktionen mit den Perioden $\omega_1$, $\omega_2$.}

Die Zahlen $a$, $b$, $c$, $\partial$ heißen die \gesperrt{Transformationszahlen}
und $n = a\partial - bc$ der \gesperrt{Transformationsgrad}.

Um die Form dieser Darstellung deutlicher zu übersehen,
wollen wir die Bedingungen aufsuchen, unter denen $T'(u, \omega_1, \omega_2)$
eine $\Theta$-Funktion der $m$ten Ordnung $\Theta(u,\omega)$ wird (§ 20). Wir
nehmen $b'_1$, $b'_2$ und folglich auch $b_1$, $b_2$ als ganze Zahlen, so daß
$b_1$, $b_2$, $b'_1$, $b'_2$ durch $g_1$, $g_2$, $g'_1$, $g'_2$ ersetzt werden können. Es
ist dann
\[
\omega_1 = 1, \quad \omega_2 = \omega, \quad a_1 = 0, \quad a_2 = m
\]
%CONT>>
